% =====================================================================
%  MIC Summer-of-Code  --  Weeks 1.5 to 3
%  Image Denoising  +  Image Reconstruction (inverse problems, CT, MRI)
%
%  Self-contained handbook chapter.  Compile with pdflatex (twice, so the
%  table of contents resolves) or upload to Overleaf and hit Recompile.
%      pdflatex MIC_Denoising_and_Reconstruction.tex
%      pdflatex MIC_Denoising_and_Reconstruction.tex
%
%  Built from Prof. Suyash P. Awate's CS-736 decks (ImageModeling,
%  ImageDenoising, ImageReconstruction, XrayCT, MRI), with attribution,
%  for the use of mentees on this project.
% =====================================================================
\documentclass[11pt]{article}

\usepackage[a4paper,margin=1in]{geometry}
\usepackage{amsmath,amssymb,amsthm,mathtools,bm}
\usepackage{enumitem}
\usepackage{booktabs}
\usepackage{xcolor}
\usepackage{microtype}
\usepackage[most]{tcolorbox}
\usepackage{titlesec}
\usepackage[colorlinks=true,linkcolor=blue!45!black,urlcolor=blue!55!black,
            citecolor=blue!45!black]{hyperref}

% ---- light styling so the document does not feel like a wall of text ----
\titleformat{\section}{\Large\bfseries\sffamily\color{blue!40!black}}{\thesection}{0.6em}{}
\titleformat{\subsection}{\large\bfseries\sffamily\color{blue!32!black}}{\thesubsection}{0.5em}{}
\titleformat{\subsubsection}{\normalsize\bfseries\sffamily}{\thesubsubsection}{0.4em}{}

\newtcolorbox{keybox}[1][]{colback=blue!4!white,colframe=blue!40!black,
  fonttitle=\bfseries\sffamily,title=#1,boxrule=0.6pt,arc=2pt,left=4pt,right=4pt,top=3pt,bottom=3pt}
\newtcolorbox{notebox}[1][Aside]{colback=black!3!white,colframe=black!45!white,
  fonttitle=\bfseries\sffamily,title=#1,boxrule=0.5pt,arc=2pt,left=4pt,right=4pt,top=3pt,bottom=3pt}

% ---- handy macros ----
\newcommand{\R}{\mathbb{R}}
\newcommand{\E}{\mathbb{E}}
\newcommand{\xhat}{\hat{x}}
\newcommand{\norm}[1]{\left\lVert #1 \right\rVert}
\newcommand{\abs}[1]{\left\lvert #1 \right\rvert}
\newcommand{\inner}[2]{\langle #1,\, #2 \rangle}
\newcommand{\deck}[1]{\textsf{\small[#1]}}
\DeclareMathOperator*{\argmin}{arg\,min}
\DeclareMathOperator*{\argmax}{arg\,max}

\title{\textbf{Image Denoising and Image Reconstruction}\\[2pt]
\large A working handbook for inverse problems, with CT and MRI physics\\[6pt]
\normalsize Medical Image Computing -- Summer of Code -- Weeks 1.5 to 3}
\author{}
\date{}

\begin{document}
\maketitle
\vspace{-2.2em}

\begin{keybox}[How to read this chapter]
This chapter is the main theory reading for this block. It is long on purpose:
you should be able to learn the whole topic from it without chasing ten
browser tabs. Read it once, slowly, start to finish. Keep the slide decks open
next to you --- wherever you see a tag like \deck{ImageDenoising, slide 14},
that is the exact slide the paragraph is explaining, so you can look at the
picture while you read the words. Do the three notebooks in \texttt{../notebooks/}
as you go; they turn every idea here into a few lines of code you can run.

One promise: nothing here is hand-waved past the point of being useful. Where a
result needs a derivation, you get the derivation. Where intuition is enough,
you get the intuition and a pointer to the depth.
\end{keybox}

\tableofcontents
\bigskip

% =====================================================================
\section{The one idea behind the whole block}
% =====================================================================

Everything in these one and a half weeks is a single sentence wearing
different costumes:

\begin{keybox}[The recurring object]
We never observe the thing we want. We observe a \emph{measurement} $y$ that is
some known transformation of the true image $x$, corrupted by noise. We recover
$x$ by asking: \emph{which $x$ best explains the data $y$, while also looking
like a plausible image?} In symbols,
\[
\xhat \;=\; \argmax_{x}\; \underbrace{p(y \mid x)}_{\text{likelihood: fits the data}}
\;\cdot\; \underbrace{p(x)}_{\text{prior: looks like an image}} .
\]
\end{keybox}

That is Bayes' theorem, and it is the spine of denoising, of CT reconstruction,
and of MRI reconstruction. The only thing that changes between those three is
the transformation that connects $x$ to $y$:

\begin{center}
\begin{tabular}{@{}llll@{}}
\toprule
Problem & True image $x$ & Measurement $y$ & Forward map \\
\midrule
Denoising      & clean image            & noisy image          & identity ($y = x + n$) \\
CT             & attenuation map        & projections (sinogram) & Radon transform \\
MRI            & magnetisation image    & $k$-space samples    & (under-sampled) Fourier \\
\bottomrule
\end{tabular}
\end{center}

So denoising is just the simplest inverse problem --- the one where the forward
map is ``do nothing''. Master it first, and CT and MRI become ``the same thing,
but with a more interesting map in the middle''. That is why we spend the first
half of this chapter on denoising and only then turn to reconstruction.

% =====================================================================
\part*{Part I -- Image Denoising}
\addcontentsline{toc}{section}{PART I -- Image Denoising}
% =====================================================================

% ---------------------------------------------------------------------
\section{What denoising is, and what noise looks like}
% ---------------------------------------------------------------------

Real scans are noisy. An MRI brain slice has a faint grainy texture; a low-dose
CT chest scan is speckled with mottle; a fetal ultrasound looks like it was
shot through frosted glass \deck{ImageDenoising, slides 3--5}. In every case the
scanner handed us $y$, and somewhere underneath is the clean $x$ we actually
want to look at or measure.

Recall from Week~1 that an image is just a long vector of numbers --- one
intensity per pixel (2D) or voxel (3D). Denoising treats each pixel's value as
a \textbf{random variable}: the scanner did not give us a fixed number, it gave
us a sample from a distribution centred (roughly) on the true intensity. Our
job is to undo the randomness.

A useful single number for ``how clean'' a measurement is, is the
\textbf{signal-to-noise ratio} \deck{ImageDenoising, slide 9}:
\[
\mathrm{SNR}(X) \;=\; \frac{\E[X]}{\mathrm{SD}(X)} ,
\]
the mean of the random variable divided by its standard deviation. Big SNR
means the signal towers over the wobble.

\begin{notebox}[i.i.d.\ versus correlated noise]
The simplest noise is \emph{independent and identically distributed} (i.i.d.):
each pixel is perturbed on its own, with the same distribution, independent of
its neighbours \deck{ImageDenoising, slide 6}. If you take an i.i.d.\ noise
image and blur it with a low-pass filter, neighbouring pixels become linked and
you get \emph{correlated} noise \deck{ImageDenoising, slides 7--8}. Most of the
classical theory assumes i.i.d.\ noise; real scanners sometimes do not oblige,
which is part of why denoising stays hard.
\end{notebox}

% ---------------------------------------------------------------------
\section{The four noise models you must know}
% ---------------------------------------------------------------------

Different physics produces different noise. Picking the right model is half the
battle, because the model \emph{is} the likelihood $p(y\mid x)$ in Bayes' rule.

\subsection{Gaussian -- the default}

The Gaussian (normal) distribution is everywhere \deck{ImageDenoising, slide 11},
for three honest reasons:
\begin{enumerate}[nosep,leftmargin=*]
  \item \textbf{Central limit theorem.} Add up many small independent effects
  and the sum looks Gaussian, almost regardless of what the small effects were.
  \item \textbf{Thermal (Johnson) noise} in the scanner's electronics is
  Gaussian; it is unavoidable above absolute zero.
  \item \textbf{Brownian motion / random walks} are Gaussian in the limit.
\end{enumerate}
Its density is
\[
p(x) \;=\; \frac{1}{\sqrt{2\pi}\,\sigma}\,
           \exp\!\left(-\frac{(x-\mu)^2}{2\sigma^2}\right),
\qquad \mathrm{SNR} = \mu/\sigma .
\]
For an additive Gaussian noise model we write $y_i = x_i + n_i$ with
$n_i \sim \mathcal{N}(0,\sigma^2)$ independently per pixel. This is the model
that makes the negative log-likelihood a clean squared error,
$-\log p(y\mid x) = \tfrac{1}{2\sigma^2}\norm{y-x}_2^2 + \text{const}$, which is
why $\ell^2$ data terms are everywhere.

\subsection{Poisson -- counting photons}

Whenever the measurement is made by \emph{counting photons} hitting a detector
--- photography, light microscopy, infrared, X-ray, CT --- the right model is
Poisson \deck{ImageDenoising, slide 14}. If on average $\lambda$ photons arrive
in the measurement window, the probability of counting exactly $k$ is
\[
\Pr(X=k) \;=\; \frac{\lambda^{k} e^{-\lambda}}{k!},
\qquad \E[X] = \mathrm{Var}(X) = \lambda,
\qquad \mathrm{SNR} = \frac{\lambda}{\sqrt{\lambda}} = \sqrt{\lambda}.
\]
Two consequences matter enormously in practice. First, the noise is
\emph{signal-dependent}: bright regions (large $\lambda$) are noisier in
absolute terms but cleaner in relative terms (SNR grows like $\sqrt\lambda$).
This is the entire argument for higher X-ray dose giving cleaner images --- and
the entire ethical problem, because dose is radiation. Second, for large
$\lambda$ the Poisson PMF starts to look like a Gaussian bell
\deck{ImageDenoising, slide 16}, which is why people often get away with a
Gaussian approximation in bright images.

\subsection{Compound Poisson -- real CT}

Practical CT uses a \emph{polychromatic} (broad-spectrum) X-ray beam, so the
count at a detector is a sum of $N$ contributions where $N$ itself is Poisson:
a \textbf{compound Poisson} variable \deck{ImageDenoising, slide 17}. It too can
be approximated by a Gaussian for many purposes, but knowing the real model
matters when you push to low dose.

\subsection{Rician -- magnitude MRI}

MRI is the subtle one. The scanner measures a \emph{complex} signal --- a
two-vector $[a,b]$ at each voxel (real and imaginary channels). Each channel
carries i.i.d.\ zero-mean Gaussian noise. But the image we look at is the
\emph{magnitude} $R=\sqrt{a^2+b^2}$, and the magnitude of a Gaussian two-vector
is \textbf{Rician} distributed \deck{ImageDenoising, slide 18}:
\[
p(x \mid \nu,\sigma) \;=\;
   \frac{x}{\sigma^2}\,
   \exp\!\left(-\frac{x^2+\nu^2}{2\sigma^2}\right)
   I_0\!\left(\frac{\nu x}{\sigma^2}\right), \qquad x \ge 0,
\]
where $\nu$ is the true (noise-free) magnitude, $\sigma$ is the per-channel
Gaussian standard deviation, and $I_0$ is the modified Bessel function of the
first kind, order $0$. Two facts to carry with you:
\begin{itemize}[nosep,leftmargin=*]
  \item $\nu$ is \emph{not} $\E[X]$ and $\sigma$ is \emph{not} $\mathrm{SD}(X)$
  --- the magnitude operation shifts and skews everything, which is exactly why
  noisy MRI never goes fully black even in air.
  \item When the true signal is zero ($\nu=0$), the Rician collapses to the
  \textbf{Rayleigh} distribution \deck{ImageDenoising, slide 20}. That is the
  noise floor you see in the background of an MR image.
\end{itemize}

\begin{notebox}[Where the Rician comes from (and a reading note)]
The derivation \deck{ImageDenoising, slides 21--22} writes the cumulative
probability $\Pr(R\le x)=\Pr(z_1^2+z_2^2\le x^2)$ as a 2D Gaussian integral
over the disc of radius $x$, switches to polar coordinates, differentiates with
the Leibniz rule to get the density, and recognises the angular integral
$\frac{1}{2\pi}\int_0^{2\pi}\exp(\nu x\cos\theta')\,d\theta' = I_0(\nu x)$ as
the very definition of $I_0$. If you open those two slides in the extracted
text you will see mangled symbols --- the PDF stored the equations as special
math glyphs that do not survive copy-paste. \textbf{Trust the slide images, not
the extracted text, on heavy-equation slides.} The clean statement is the
boxed density above.
\end{notebox}

\subsection{Ultrasound speckle}

Ultrasound noise is genuinely messy --- it depends on the device, the wave
coherence, and the tissue's scattering density, and people model it with
Rician, Rayleigh, or K-distributions \deck{ImageDenoising, slide 23}. A common
working model is \emph{multiplicative speckle}
\deck{ImageDenoising, slide 45}:
\[
Y = X + \sqrt{X}\,Z, \qquad Z\sim\mathcal{N}(0,\sigma^2),
\]
so that the noise grows with the signal (its standard deviation scales like
$\sqrt{X}$). You recover the likelihood $p(y\mid x)$ by the change-of-variables
formula for densities.

\subsection{Two facts that pay off later}

\textbf{Averaging helps by $\sqrt{N}$.} If you acquire the same voxel $N$ times
with i.i.d.\ noise and average, the signal mean grows like $N$ but the standard
deviation only like $\sqrt N$, so SNR improves by $\sqrt N$
\deck{ImageDenoising, slide 24}. This is why slow MRI sequences average repeats,
and why ``just scan longer'' is the lazy way to a cleaner image.

\textbf{Variance stabilisation.} It is sometimes convenient to turn
signal-dependent noise into roughly constant-variance noise so you can pretend
it is Gaussian. For Poisson data the \textbf{Anscombe transform}
\[
f(X) = 2\sqrt{X + \tfrac{3}{8}}
\]
maps a Poisson variable to something with variance $\approx 1$
\deck{ImageDenoising, slide 30}. Denoise in the transformed domain, then invert.
It is a trick, not a religion --- the course prefers modelling the real noise
--- but it is good to recognise.

% ---------------------------------------------------------------------
\section{Priors: teaching the computer what an image looks like}
% ---------------------------------------------------------------------

The likelihood alone cannot denoise. If we only maximise $p(y\mid x)$ with a
per-pixel Gaussian, the best answer is $\xhat = y$: ``the cleanest explanation
of the noisy image is the noisy image itself''. We need a \textbf{prior}
$p(x)$ that says some images are more plausible than others. Three physical
beliefs drive almost every image prior \deck{ImageModeling, slide 16}:
\begin{enumerate}[nosep,leftmargin=*]
  \item intensities usually vary \emph{smoothly} across space;
  \item big jumps happen, but only at \emph{object boundaries};
  \item the number of objects is far smaller than the number of pixels.
\end{enumerate}

\subsection{Random fields on a grid}

To turn ``neighbouring pixels are similar'' into probability, we put a
\textbf{random field} on the pixel lattice \deck{ImageModeling, slide 19}: a
collection of random variables $\{X_i\}$, one per site $i$, together with a
\textbf{neighbourhood system} $N=\{N_i\}$ that lists which sites are neighbours
\deck{ImageModeling, slide 20}. In 2D the usual choices are the 4-neighbour
(up/down/left/right) or 8-neighbour systems; in 3D, 6/18/26 neighbours. A
\textbf{clique} is a set of sites that are all mutually neighbours --- single
pixels, neighbouring pairs, and so on \deck{ImageModeling, slide 21}.

\subsection{Markov Random Fields}

A \textbf{Markov Random Field} (MRF) is a random field with a local-memory
property \deck{ImageModeling, slide 22}:
\[
p(X_i \mid X_{S\setminus\{i\}}) \;=\; p(X_i \mid X_{N_i}).
\]
In words: \emph{a pixel, given its neighbours, forgets about the rest of the
image.} This is the spatial version of the Markov property. A common
misconception, worth nailing now: this does \emph{not} say that far-apart pixels
are independent. They are only \emph{conditionally} independent given the
pixels between them; information still propagates across the whole image,
neighbour to neighbour.

\subsection{Gibbs fields, energy, and temperature}

MRFs are defined by their local conditionals, which is awkward --- we would
rather write down the joint distribution of the whole image. The
\textbf{Hammersley--Clifford theorem} (1971) is what lets us
\deck{ImageModeling, slides 24, 35--43}:

\begin{keybox}[Hammersley--Clifford]
A random field is an MRF with respect to a neighbourhood system \emph{if and
only if} it is a \textbf{Gibbs random field} with respect to the same system
(assuming all configurations have positive probability). Equivalently, its
joint distribution is a Gibbs distribution
\[
p(x) \;=\; \frac{1}{Z}\,\exp\!\left(-\frac{U(x)}{T}\right),
\qquad
U(x) = \sum_{c \in \mathcal{C}} V_c(x_c).
\]
\end{keybox}

Here $U$ is the \textbf{energy}, a sum of \textbf{clique potentials} $V_c$ that
you design to encode your beliefs; $Z$ (the \textbf{partition function}) is the
normaliser that makes the probabilities sum to one; and $T$ is a
\textbf{temperature} \deck{ImageModeling, slides 25--29}. Low energy means high
probability, so a ``plausible'' image is a \emph{low-energy} image. The
temperature controls how peaked the distribution is: as $T\to 0$ all the mass
piles onto the lowest-energy images, and as $T\to\infty$ every image becomes
equally likely (a uniform distribution). By default we take $T=1$.

The partition function $Z$ is the villain of the story: summing over every
possible image to compute it is hopeless. The saving grace is that for finding
the \emph{most probable} image we maximise $p(x)$ over $x$, and $Z$ does not
depend on $x$, so it drops out. We almost never need $Z$ itself.

\subsection{Why a naive prior blurs edges, and how to fix it}

Now the design question: what should the clique potential be? Take the simplest
useful case --- a pixel $x$ with neighbours $y_i = x + r_i$, where $r_i$ is the
deviation of neighbour $i$ \deck{ImageModeling, slide 45}. Model the prior
energy as a sum of penalties on deviations,
\[
E(x) = \sum_i g(y_i - x),
\]
and remember that the most-probable $x$ minimises this energy
\deck{ImageModeling, slide 46}. Choose $g$ even, and write $g(u)=H(|u|^2)$. The
stationarity condition $E'(x)=0$ gives a beautiful interpretation: the update
on $x$ is a \emph{weighted average} of its neighbours, with weights
$h(u) = g'(u)/(2u)$ acting as an interaction strength \deck{ImageModeling,
slides 47--48}. The force that neighbour $y_i$ exerts on $x$ is proportional to
$(y_i-x)\,h(y_i-x)$.

This single picture explains the central tension of image priors:

\begin{itemize}[nosep,leftmargin=*]
  \item \textbf{Quadratic penalty} $g(u)=u^2$ gives $h(u)=1$
  \deck{ImageModeling, slide 49}. This is a Gaussian prior. The force
  $u\cdot h(u)=u$ keeps \emph{growing} with the deviation, so a real edge (a
  big $u$) pulls just as hard as noise --- the model fights every jump and
  \textbf{blurs edges}. It also makes discontinuities ``extremely unlikely'',
  which is exactly wrong for medical images full of boundaries.
  \item \textbf{Edge-preserving penalties} are designed so the force
  \emph{saturates}: as $u\to\infty$, $h(u)\to 0$ and $u\,h(u)$ stays bounded
  \deck{ImageModeling, slides 50--56}. Then a large deviation is recognised as
  ``probably a real edge, leave it alone''. The Huber penalty (quadratic near
  zero, linear far out) and the generalised-Gaussian family are the workhorses.
\end{itemize}

\begin{keybox}[The edge-preservation rule of thumb]
A good image prior should penalise \emph{small} deviations like a quadratic (to
kill noise in flat regions) and penalise \emph{large} deviations only gently
--- at most linearly (to leave edges intact). Quadratic-everywhere = smooth
everywhere = blurred. This is the same reason $\ell^1$ beats $\ell^2$, which we
return to in Part~II.
\end{keybox}

% ---------------------------------------------------------------------
\section{Bayesian denoising as MAP estimation}
% ---------------------------------------------------------------------

Now assemble the pieces. The \textbf{posterior} is likelihood times prior,
normalised \deck{ImageDenoising, slide 25}. The \textbf{maximum a posteriori}
(MAP) estimate is the image that maximises it:
\[
\xhat
  = \argmax_x\; p(y\mid x)\,p(x)
  = \argmin_x\; \Big[\,\underbrace{-\log p(y\mid x)}_{\text{data fidelity}}
        \;\underbrace{-\log p(x)}_{\text{regularisation}}\,\Big].
\]
Taking the negative logarithm turns the product into a sum and the search for a
peak into the more familiar search for a valley. The two terms have crisp jobs:
the \textbf{data-fidelity} term forces agreement with the measurement; the
\textbf{regularisation} term forces the answer to look like a plausible image.

For an i.i.d.\ Gaussian noise model the likelihood factorises across pixels
\deck{ImageDenoising, slide 28}, $p(y\mid x)=\prod_i \mathcal{N}(y_i\mid
x_i,\sigma^2)$, and the data term becomes $\tfrac{1}{2\sigma^2}\norm{y-x}_2^2$.
With a Gibbs prior of energy $U$, the whole objective is
\[
\xhat = \argmin_x\; \frac{1}{2\sigma^2}\norm{y-x}_2^2 + U(x).
\]

\begin{keybox}[Prior $=$ regulariser: the dictionary]
A negative-log Gibbs prior \emph{is} a regulariser. This is the single most
useful equivalence in the course.
\[
\text{Gaussian prior} \;\leftrightarrow\; \ell^2\ \text{(Tikhonov) penalty},
\qquad
\text{Laplace prior} \;\leftrightarrow\; \ell^1\ \text{(TV/Lasso) penalty}.
\]
Whenever you see a regularised least-squares problem, you are looking at a MAP
estimate in disguise, and vice versa.
\end{keybox}

% ---------------------------------------------------------------------
\section{Actually solving it: ICM, gradient descent, annealing}
% ---------------------------------------------------------------------

The MAP objective is usually too big and too coupled to solve in one shot, so
we optimise iteratively.

\paragraph{Iterated Conditional Modes (ICM).} The simplest idea: fix every
pixel except one, and move that one pixel to the value that maximises its
\emph{local} conditional posterior; then sweep over all pixels and repeat
\deck{ImageDenoising, slides 37--39}. The local objective is just
(local likelihood) $\times$ (local conditional prior), which thanks to the MRF
property only involves the pixel's neighbours. Stop when a full sweep barely
changes the image (for example when the RMS change drops below $1\%$). ICM is
easy and monotone, but it is greedy and can get stuck.

\paragraph{Parallel gradient updates.} ICM visits pixels one at a time, and the
visiting order can leave artefacts \deck{ImageDenoising, slide 40}. A practical
alternative is to step every pixel towards its local mode at once (gradient
ascent on the posterior), and accept the step if the posterior improved. For
the Rician and speckle likelihoods the per-voxel derivatives involve the Bessel
functions $I_0$ and $I_1$ \deck{ImageDenoising, slides 43, 46}; the algebra is
fiddly but the principle is identical.

\paragraph{Balancing the two terms.} In practice you weight the prior and the
likelihood against each other with parameters $\beta$ (prior strength) and
$\alpha$ (likelihood strength) \deck{ImageDenoising, slides 47--52}. The two
extremes are illuminating: $\beta=0$ removes the prior entirely and you fall
back to the \textbf{maximum-likelihood} estimate (pure data fit, noisy);
$\alpha=0$ makes the likelihood flat and the answer is driven only by the prior
(over-smoothed). The art is in between, and that knob is the single most
important thing you will tune in the notebooks.

\paragraph{Simulated annealing (escaping local minima).} When the energy is
non-convex --- especially for discrete labelling --- gradient methods get
trapped. \textbf{Simulated annealing} \deck{ImageModeling, slides 30--34}
borrows from metallurgy: propose a random move; always accept it if it lowers
the energy; accept it \emph{even if it raises} the energy with probability
$p(y)/p(x)$, which at high temperature is generous. Start hot (wander freely,
ignore where you started), then cool slowly so the walker settles into a deep
valley. Cool too fast and you freeze into the nearest dent; the theory needs a
painfully slow schedule like $T(n)=d/\log n$ for a global-optimum guarantee.

% ---------------------------------------------------------------------
\section{The modern face of the smoothness prior: total variation}
% ---------------------------------------------------------------------

The edge-preserving idea has a famous continuous cousin. \textbf{Total
variation} (TV) measures the total amount of intensity change in an image,
\[
\mathrm{TV}(x) = \int \norm{\nabla x}\, ,
\]
i.e.\ the $\ell^1$ norm of the gradient magnitude. The
Rudin--Osher--Fatemi (ROF) model denoises by solving
\[
\xhat = \argmin_x\; \tfrac12\norm{x-y}_2^2 + \lambda\,\mathrm{TV}(x).
\]
Because TV uses an $\ell^1$ norm of the gradient, it is happy to leave a few
large jumps (edges) in exchange for flattening everything else --- precisely
the edge-preserving behaviour we engineered by hand above, now stated in one
line. Its tell-tale artefact is the ``staircase'' effect: smooth ramps turn
into little flat steps. TV is non-differentiable at zero, so it is solved with
proximal/primal--dual methods rather than plain gradient descent (see
\href{https://angms.science/doc/TV/TV.pdf}{A.\,M.-C.\ So's ROF notes} and the
\href{https://sangillee.com/2026-01-06-total-variation-minimization/}{Chambolle--Pock walkthrough}).
You will implement a simple version in Notebook~3.

% =====================================================================
\part*{Part II -- Inverse Problems and Bayesian Reconstruction}
\addcontentsline{toc}{section}{PART II -- Inverse Problems and Reconstruction}
% =====================================================================

% ---------------------------------------------------------------------
\section{From denoising to reconstruction: the forward model}
% ---------------------------------------------------------------------

In denoising the measurement lived in the same space as the image. In
\textbf{reconstruction} the scanner first \emph{transforms} the image and then
adds noise \deck{ImageReconstruction, slide 2}:
\[
y = G x + n,
\]
where $G$ is the \textbf{system operator} (a big matrix, conceptually) that
models the physics of acquisition, and $n$ is noise. The image $x$ is the
physical quantity of interest --- attenuation coefficients in CT, magnetisation
in MRI. The prior on $x$ is exactly the same machinery as denoising (an MRF, a
TV prior, a dictionary). All that is new is $G$.

With i.i.d.\ Gaussian noise, the negative log-likelihood is again a squared
error, now \emph{through} $G$ \deck{ImageReconstruction, slide 3}:
\[
-\log p(y\mid x) = \tfrac{1}{2\sigma^2}\norm{Gx - y}_2^2 + \text{const}.
\]

% ---------------------------------------------------------------------
\section{Why reconstruction is hard: ill-posedness}
% ---------------------------------------------------------------------

If $G$ were nicely invertible we would just compute $x = G^{-1}y$ and go home.
It never is. Reconstruction problems are \textbf{ill-posed}
\deck{ImageModeling, slides 69--70}:
\begin{itemize}[nosep,leftmargin=*]
  \item $G$ may have a non-trivial \textbf{null space} --- many different images
  produce \emph{exactly} the same measurements, so the data alone cannot pick
  one (under-determined).
  \item Even when a unique least-squares solution exists, $G$ can be
  \textbf{ill-conditioned}: tiny measurement noise gets amplified into huge
  swings in $x$. The objective has long, flat ``ridges'' along which many very
  different images fit almost equally well --- the origin of the name
  \emph{ridge} regression \deck{ImageModeling, slide 70}.
\end{itemize}
The cure is the prior. By adding $-\log p(x)$ we tilt the flat ridge so that one
plausible image becomes the clear winner. This is regularisation, and it is not
optional --- it is what makes the problem solvable at all.

% ---------------------------------------------------------------------
\section{The MAP reconstruction objective}
% ---------------------------------------------------------------------

Putting likelihood and prior together gives the objective that the rest of the
course keeps reusing:
\[
\xhat \;=\; \argmin_x\; \frac{1}{2\sigma^2}\norm{Gx-y}_2^2 \;+\; \lambda\,R(x).
\]
Two classic choices of regulariser $R$, both seen already as priors:

\paragraph{Tikhonov / ridge ($\ell^2$).} $R(x)=\norm{\Gamma x}_2^2$
\deck{ImageModeling, slides 70--74}. The objective stays quadratic, so it has a
closed-form solution
\[
\xhat = (G^{\!\top}G + \Gamma^{\!\top}\Gamma)^{-1} G^{\!\top} y,
\]
and a clean Bayesian reading: it is the MAP estimate under a Gaussian
likelihood and a zero-mean Gaussian prior with covariance
$(\Gamma^{\!\top}\Gamma)^{-1}$. Easy to solve, but it smooths edges (it is the
quadratic prior again).

\paragraph{TV / Lasso ($\ell^1$).} $R(x)=\norm{x}_1$ or $\mathrm{TV}(x)$
\deck{ImageModeling, slides 75--84}. The Bayesian reading is a Laplace prior.
The payoff is \textbf{sparsity}: $\ell^1$ drives small coefficients to exactly
zero rather than merely shrinking them. Geometrically, the $\ell^1$ ball has
sharp corners on the axes, and the solution of a constrained fit tends to land
on a corner --- where coordinates are zero \deck{ImageModeling, slides 80--81}.
That is why $\ell^1$ recovers sparse signals and edge maps where $\ell^2$ gives
a blurry compromise. The ideal sparsity penalty is the $\ell^0$ ``norm'' (count
of non-zeros), but it is NP-hard, and $\ell^1$ is the tractable convex
relaxation that still promotes sparsity \deck{ImageModeling, slide 84}. Real
images make this concrete: the histogram of wavelet coefficients of natural
images is sharply peaked at zero and heavy-tailed --- decidedly non-Gaussian,
exactly the shape an $\ell^1$/Laplace prior expects \deck{ImageModeling,
slide 85}.

% ---------------------------------------------------------------------
\section{The optimisation toolbox}
% ---------------------------------------------------------------------

Reconstruction objectives are minimised numerically, so you need a small
optimisation toolkit \deck{ImageReconstruction, slides 4--38}.

\subsection{The one gradient you must be able to derive}

Everything starts from the gradient of the data term. With $f(x)=\norm{Gx-y}_2^2
= (Gx-y)^{\!\top}(Gx-y)$, set $w=Gx-y$ and use the chain rule
\deck{ImageReconstruction, slide 8}:
\[
\frac{\partial f}{\partial w} = 2w^{\!\top},
\qquad
\frac{\partial w}{\partial x} = G,
\qquad\Rightarrow\qquad
\frac{\partial f}{\partial x} = 2(Gx-y)^{\!\top}G,
\]
and the gradient is the transpose of that derivative:
\[
\boxed{\;\nabla_x \norm{Gx-y}_2^2 = 2\,G^{\!\top}(Gx-y).\;}
\]
Notice the two operators: $G$ takes you image$\,\to\,$data (the forward map),
and $G^{\!\top}$ takes you data$\,\to\,$image (the \emph{adjoint}, or
back-projection). Every iterative reconstruction is a loop of ``apply $G$, see
the mismatch, push it back with $G^{\!\top}$''.

\begin{notebox}[Row vectors, gradients, and a notation gotcha]
The derivative of a scalar function is naturally a \emph{row} vector; the
\emph{gradient} is its transpose, a column vector the same shape as $x$
\deck{ImageReconstruction, slides 5--6}. For a vector-valued function the
derivative is the \textbf{Jacobian} (rows = outputs, columns = inputs). Keeping
shapes straight is most of the battle in matrix calculus; the
\href{https://www.math.uwaterloo.ca/~hwolkowi/matrixcookbook.pdf}{Matrix
Cookbook} is the reference to keep open.
\end{notebox}

\subsection{Gradient descent and step sizes}

Gradient descent just walks downhill \deck{ImageReconstruction, slides 18--22}:
\[
x_{k+1} = x_k - s\,\nabla f(x_k).
\]
The whole game is the step size $s$. A robust, derivative-free rule
\deck{ImageReconstruction, slide 22}: propose a step; if the objective went
down, accept it and grow $s$ a little (say $\times 1.1$) for next time; if it
went up, halve $s$ and try again. This adapts to the local landscape without an
expensive line search. Formal line searches (golden-section, Newton,
quasi-Newton) exist but are often skipped when an approximate step is good
enough \deck{ImageReconstruction, slide 19}.

\subsection{Second-order and least-squares methods}

\textbf{Newton's method} uses curvature: expand $f$ to second order and jump to
the minimum of that quadratic, giving the update $H\,\Delta x = -\nabla f$ with
$H$ the \textbf{Hessian} \deck{ImageReconstruction, slides 26--32}. It converges
fast near the optimum but needs second derivatives. For least-squares problems
$S(\beta)=\sum_i r_i(\beta)^2$, the \textbf{Gauss--Newton} method approximates
the Hessian by $J^{\!\top}J$ (just first derivatives of the residuals), and
\textbf{Levenberg--Marquardt} damps it,
\[
\Delta = -s\,(J^{\!\top}J + \lambda I)^{-1} J^{\!\top} r,
\]
interpolating between Gauss--Newton (small $\lambda$) and plain gradient descent
(large $\lambda$) \deck{ImageReconstruction, slides 35--38}. LM is the default
for nonlinear least squares.

\subsection{Constraints: projected gradient and POCS}

Often $x$ must obey constraints --- attenuation coefficients cannot be negative,
for instance. \textbf{Projected gradient descent} takes an ordinary step and
then projects back onto the constraint set; for convex sets it behaves like
ordinary gradient descent \deck{ImageReconstruction, slide 25}. Enforcing
several convex constraints by cycling through their projections is the method of
\textbf{Projections Onto Convex Sets} (POCS), which we meet again in CT.

\subsection{Complex derivatives (you will need them for MRI)}

MRI images are complex-valued, and differentiating with respect to a complex
variable is genuinely different from the real case, because the limit must hold
as the step approaches zero \emph{from every direction in the plane}
\deck{ImageReconstruction, slides 9--14}. Forcing the real-axis and
imaginary-axis limits to agree gives the \textbf{Cauchy--Riemann equations}
$u_x=v_y,\ u_y=-v_x$. The practical escape hatch used in reconstruction code is
the \emph{complex-as-doubled-real} trick: represent $w=x+iy$ as the real vector
$[x,y]^{\!\top}$ and optimise there. Then for $f(w)=w^{*}w=\abs{w}^2$ you get the
tidy result $\nabla_w f = 2w$, and gradient descent on complex images proceeds
just like the real case.

% ---------------------------------------------------------------------
\section{Algebraic reconstruction (ART)}
% ---------------------------------------------------------------------

There is an alternative to optimising a global objective: enforce one
measurement equation at a time. \textbf{ART} (Algebraic Reconstruction
Technique), the Kaczmarz method by another name, writes the system as
$Ax=b$ and repeatedly nudges the current image to satisfy each row
\deck{XrayCT, slides 85--92}:
\[
x \;\leftarrow\; x + \lambda_k\,\frac{b_i - \inner{a_i}{x}}{\norm{a_i}^2}\,a_i,
\]
where $a_i^{\!\top}$ is the $i$-th row of $A$ and $0<\lambda_k\le 1$ is a
relaxation parameter. Geometrically, each row defines a hyperplane of consistent
images, and the update \emph{projects} onto it --- POCS again. ART shines
exactly when the clean Fourier inversion (next part) fails: few view angles,
metal artefacts, very low dose.

% =====================================================================
\section{CT and MRI as forward operators (a preview)}
% =====================================================================

The two big medical-imaging inverse problems are just our reconstruction
template with a specific operator $G$ in the middle. We preview them in one
paragraph each; the full physics and modelling intuition --- Beer's law, the
Radon transform, filtered back-projection, the Bloch equations, $k$-space --- get
their own dedicated chapter in \texttt{week\_3} (CT \& MRI fundamentals), with a
companion notebook.

\paragraph{CT: $G$ is the Radon transform.} An X-ray measures the line integral
of tissue attenuation along its path. Collect such integrals over many angles
and you have the \textbf{Radon transform} of the attenuation image --- the
sinogram. Reconstruction inverts it: analytically by \textbf{filtered
back-projection}, or, when data is noisy or sparse, by the same MAP objective
$\min_x \tfrac12\norm{Gx-y}_2^2+\lambda R(x)$ with $G$ the Radon operator and
$G^{\!\top}$ the back-projector.

\paragraph{MRI: $G$ is the Fourier transform.} An MRI scanner measures the 2D
\textbf{Fourier transform} of the slice (``$k$-space''), one sample at a time.
Reconstruction is an inverse FFT when $k$-space is full; when it is
under-sampled (to scan faster) it becomes an ill-posed inverse problem cured by
a prior --- compressed-sensing MRI. Same template, $G$ a (subsampled) Fourier
operator.

\begin{keybox}[Where the physics lives]
Everything about \emph{how} CT and MRI actually form $G$ and $y$ --- photon
counting, attenuation and Hounsfield units, the Fourier Slice Theorem, spins,
$T_1/T_2$ relaxation, RF excitation and spatial encoding --- is in the
\texttt{week\_3} handbook (\texttt{MIC\_CT\_and\_MRI\_Fundamentals.tex}). Here we
only needed to know that CT and MRI are \emph{linear forward operators}, so the
Bayesian reconstruction of Part~II applies to them unchanged.
\end{keybox}

% =====================================================================
\section{The picture to carry forward}
% =====================================================================

Step back and the three modalities are one method:
\[
\xhat = \argmin_x\; \tfrac{1}{2\sigma^2}\norm{Gx-y}_2^2 + \lambda R(x),
\qquad
\begin{cases}
G=I & \text{denoising},\\
G=\text{Radon} & \text{CT},\\
G=\text{Fourier} & \text{MRI}.
\end{cases}
\]
Choose the data term from the noise physics (Gaussian, Poisson, Rician); choose
the regulariser from what images look like (smooth with edges $\Rightarrow$
edge-preserving / TV / $\ell^1$); then descend. Denoising is the special case
that taught us the prior; CT and MRI are the same optimisation with a Radon or
Fourier operator in the middle and real physics setting $G$ and the noise. Hold
on to that and the rest of the course is variations on a theme.

\begin{keybox}[Do the notebooks]
\textbf{Notebook 1} builds the four noise models and shows their histograms,
SNR, and the Anscombe transform. \textbf{Notebook 2} implements MAP denoising
with an MRF prior (ICM and gradient descent) and lets you watch a quadratic
prior blur edges while a Huber prior keeps them. \textbf{Notebook 3} sets up the
forward model $y=Gx+n$ and reconstructs by deblurring (naive inverse vs Tikhonov
vs TV) --- the general inverse problem. (The CT Radon/FBP and MRI $k$-space
hands-on live in the \texttt{week\_3} notebook.) Run them, then change the
numbers and break things --- that is where the understanding sticks.
\end{keybox}

% =====================================================================
\section*{Resources (watch one, read one)}
\addcontentsline{toc}{section}{Resources}
% =====================================================================

\textbf{Bayesian / MRF foundations}
\begin{itemize}[nosep,leftmargin=*]
  \item Geman \& Geman (1984), \href{https://doi.org/10.1109/TPAMI.1984.4767596}{\emph{Stochastic Relaxation, Gibbs Distributions, and the Bayesian Restoration of Images}} --- the founding paper (sections I--III).
  \item C.\ Bouman, \href{https://engineering.purdue.edu/~bouman/publications/pdf/MBIP-book.pdf}{\emph{Model-Based Image Processing}} --- free book; chapters 1--2 then 7--8.
  \item Mohammad-Djafari, \href{https://www.mdpi.com/1099-4300/23/12/1673}{\emph{Regularization, Bayesian Inference, and ML Methods for Inverse Problems}} (Entropy 2021) --- tutorial covering denoising, restoration, and CT in one place.
\end{itemize}

\textbf{Noise models}
\begin{itemize}[nosep,leftmargin=*]
  \item Gudbjartsson \& Patz (1995), \href{https://onlinelibrary.wiley.com/doi/10.1002/mrm.1910340618}{\emph{The Rician Distribution of Noisy MRI Data}} --- four very readable pages.
  \item \href{https://en.wikipedia.org/wiki/Anscombe_transform}{Anscombe transform} (Wikipedia) for variance stabilisation.
\end{itemize}

\textbf{Denoising / total variation}
\begin{itemize}[nosep,leftmargin=*]
  \item Steve Brunton, \href{https://www.youtube.com/watch?v=DLfMsBKaDOY}{\emph{Sparsity and the L1 norm}} (video intuition).
  \item A.\ M.-C.\ So, \href{https://angms.science/doc/TV/TV.pdf}{\emph{Total Variation and the ROF denoising problem}} (concise derivation).
  \item \href{https://sangillee.com/2026-01-06-total-variation-minimization/}{TV minimisation with Chambolle--Pock} (worked tutorial); Buades et al.\ (2005), \href{https://ieeexplore.ieee.org/document/1467423}{\emph{A non-local algorithm for image denoising}}.
\end{itemize}

\textbf{Inverse problems / optimisation}
\begin{itemize}[nosep,leftmargin=*]
  \item \href{https://computational-imaging.de/book/06_InverseProblemsInCompimg.html}{Inverse Problems in Computational Imaging} --- course notes with the MAP derivation.
  \item Steve Brunton, \href{https://www.youtube.com/watch?v=Jl3Ftf6S5MQ}{\emph{Inverse problems}} (video); Petersen \& Pedersen, \href{https://www.math.uwaterloo.ca/~hwolkowi/matrixcookbook.pdf}{\emph{The Matrix Cookbook}}.
  \item Boyd \& Vandenberghe, \href{https://web.stanford.edu/~boyd/cvxbook/}{\emph{Convex Optimization}}, ch.\ 9 (gradient/Newton methods).
\end{itemize}

\textbf{CT and MRI.} The full physics, modelling intuition, and resource list
live in the dedicated \texttt{week\_3} handbook (CT \& MRI fundamentals).

\vfill
\begin{center}\small\itshape
Built for the MIC Summer-of-Code from Prof.\ Suyash P.\ Awate's CS-736 slide
decks, with attribution. For mentee use on this project; please do not
redistribute the bundled decks outside it.
\end{center}

\end{document}
